% Part: lambda-calculus
% Chapter: lambda-definability
% Section: fixpoints

\documentclass[../../../include/open-logic-section]{subfiles}

\begin{document}

\olfileid{lam}{ldf}{fp}

\olsection{સ્થિરબિંદુઓ}

ધારો કે આપણે ક્રમગુણિત વિધેયને પુનરાવર્તન દ્વારા એવા
પદ~$\fn{Fac}$ તરીકે વ્યાખ્યાયિત કરવા માગીએ છીએ, જેમાં નીચેનો ગુણધર્મ હોય:
\[
\fn{Fac} \ident \lambd[n][\fn{IsZero}\, n\, \num 1 (\fn{Mult}\, n (\fn{Fac}(\fn{Pred} \, n)))]
\]
એટલે $n=0$ હોય ત્યારે $n$નું ક્રમગુણિત~$1$ છે અને
અન્યથા તે $n$ ગુણ્યા $n-1$નું ક્રમગુણિત છે. અલબત્ત, આ રીતે
$\fn{Fac}$ પદ વ્યાખ્યાયિત કરી શકાતું નથી, કારણ કે જમણી બાજુએ
$\fn{Fac}$ પોતે જ આવે છે. પોતાનો જ સંદર્ભ લેતી આવી
પુનરાવર્તી વ્યાખ્યાઓ લૅમ્બડા કલનનો ભાગ નથી. દાખલા તરીકે,
\[
\fn{Mult} \ident \lambd[ab][a (\fn{Add}\, a) 0]
\]
વડે પદ વ્યાખ્યાયિત કરીએ ત્યારે જમણી બાજુએ $\fn{Add}$ જેવાં
પહેલેથી વ્યાખ્યાયિત પદો જ આવે છે. $\fn{Add}$ને તેનું વ્યાખ્યાયક
પદ મૂકીને હંમેશા દૂર કરી શકીએ છીએ. આમ $\fn{Mult}$નું
શુદ્ધ લૅમ્બડા પદ મળે; જો $\fn{Add}$ પોતે વ્યાખ્યાયિત પદો
ધરાવતું હોય (જેમ કે $\fn{Add}'$ ધરાવે છે), તો આ પ્રક્રિયા
ચાલુ રાખીને છેવટે શુદ્ધ લૅમ્બડા પદ સુધી પહોંચી શકીએ.

પરંતુ ઉપરની $\fn{Fac}$ જેવી પુનરાવર્તી વ્યાખ્યામાં એવું થતું
નથી. જમણી બાજુએ આવતા $\fn{Fac}$ના સ્થાને
$\fn{Fac}$ની પોતાની વ્યાખ્યા મૂકીએ તો મળે છે:
\begin{align*}
  \fn{Fac} & \ident \lambd[n][\fn{IsZero}\, n\, \num{1}\,
    (\fn{Mult}\, n
  ((\lambd[n][\fn{IsZero} \, n \, \num 1\, (\fn{Mult}\, n\, (\fn{Fac}
    (\fn{Pred}\, n)))]) (\fn{Pred}\, n)))]
\end{align*}
\sourcecorrection{OLLAM-057}{સ્થિર અંગ્રેજી મૂળના આ
અવેજીકરણમાં $\fn{IsZero}$ના ત્રીજા આર્ગ્યુમેન્ટ પહેલાં
લૅમ્બડા પદ બંધ થઈ જાય છે અને $\fn{Mult}$ બહાર રહી જાય છે.
અહીં એ શાખાને પદની અંદર જ રાખી છે, જેથી ઉપરની
પુનરાવર્તી સમીકરણમાં કરેલો અવેજ સાચો રહે.}
જમણી બાજુમાંથી $\fn{Fac}$ હજુ દૂર થયું નથી. આ પ્રક્રિયા
પુનરાવર્તિત કરીએ તો વ્યાખ્યા લાંબી થતી જાય છે, પરંતુ
શુદ્ધ લૅમ્બડા પદ કદી મળતું નથી. તેથી ક્રમગુણિતને
(અથવા સામાન્ય રીતે પુનરાવર્તી વિધેયોને) આ રીતે વ્યાખ્યાયિત
કરવું શક્ય નથી.

તેમ છતાં પુનરાવર્તી વ્યાખ્યા આપણને કંઈક કહે છે: જો
$f$ ક્રમગુણિત વિધેયનું નિરૂપણ કરતું પદ હોય, તો
\[
\fn{Fac}' \ident \lambd[g][\lambd[n][\fn{IsZero} \, n \, \num 1\, (\fn{Mult} \, n\, (g (\fn{Pred} n)))]]
\]
પદને $f$ પર લાગુ કરતાં મળતું $\fn{Fac}'\,f$ પણ ક્રમગુણિત
વિધેયનું નિરૂપણ કરે છે. એટલે $\fn{Fac}'$ને વિધેય સ્વીકારીને
વિધેય પરત કરતું વિધેય ગણીએ તો, $f$ ક્રમગુણિત વિધેયનું
નિરૂપણ કરતું હોય ત્યારે $\fn{Fac'}\, f$ એ જ વિધેયનું
નિરૂપણ કરે છે. જે પદ~$f$ માટે વધુમાં $\fn{Fac}' \, f
\equal[\beta] f$ હોય, તેને $\fn{Fac}'$નું
\emph{સ્થિરબિંદુ} કહે છે. વિધેયોના સ્તરે ક્રમગુણિત
$\fn{Fac}'$ હેઠળ અચળ રહે છે; આવું બીટા-સ્થિરબિંદુ
આપતું ચોક્કસ પદ~$f$ નીચે બનાવશું.
\sourcecorrection{OLLAM-059}{સ્થિર અંગ્રેજી મૂળમાં
બંને પદો એક જ ક્રમગુણિત વિધેયનું નિરૂપણ કરે છે
એ હકીકતથી પદોની $\beta$-સમતા તરફ સીધી છલાંગ છે.
સમાન વિધેયનું નિરૂપણ માત્ર $\beta$-સમતા આપતું નથી.
અહીં વિધેયનું અચળપણું અને પદની વધારાની
$\beta$-સ્થિરબિંદુ શરત અલગ દર્શાવી છે.}

લૅમ્બડા કલનમાં આપેલા પદનાં સ્થિરબિંદુઓ ગણતાં પદો છે;
તેમનો ઉપયોગ કરીને $\fn{Fac}'$ જેવા પદમાંથી ક્રમગુણિતની
વ્યાખ્યા મેળવી શકાય છે.

\begin{defn}
  \ollabel{defn:Turing-Y}
  \emph{Y-સંયોજક} આ પદ છે:
  \[
  Y \ident (\lambd[ux][x(uux)])(\lambd[ux][x(uux)]).
  \]
\end{defn}

\begin{thm}
  $Y$નો એવો ગુણધર્મ છે કે કોઈ પણ પદ~$g$ માટે
  $Yg \red g(Yg)$ છે. તેથી
  $Yg$ હંમેશા $g$નું સ્થિરબિંદુ છે.
\end{thm}

\begin{proof}
  $(\lambd[ux][x(uux)])$ને સંક્ષેપમાં $U$ કહીએ; ત્યારે
  $Y \ident UU$. હવે
  \begin{align*}
    Y g &\ident (\lambd[ux][x(uux)])U\,g \\
    &\red (\lambd[x][x(UUx)])g\\
    & \red g(UUg) \ident g(Yg).
  \end{align*}
  $g(Yg)$ અને $Yg$ બંને $g(Yg)$માં લઘુકૃત થાય છે,
  તેથી $g(Yg) \equal[\beta] Yg$; આમ $Yg$ એ $g$નું
  સ્થિરબિંદુ છે.
\end{proof}

અલબત્ત, $Yg$ લઘુકરણીય પદ હોવાથી લઘુકરણ અનંત સુધી
ચાલુ રહી શકે છે:
\begin{align*}
  Y g &\red g (Y g) \\
    &\red g (g (Y g)) \\
    &\red g(g (g (Y g)))\\
    &\ldots
\end{align*}
તેથી $Yg$ને જાણે $g$ને પોતે જ અનંત વખત લાગુ કરવાથી
મળતું પદ ગણી શકીએ. તેના પર $g$ એક વધુ વખત લાગુ
કરીએ તો, કહીએ તો, કંઈ વધારાનું થતું નથી: $Yg$ પર
અનંત વખત લાગુ થયેલા $g$ પર ફરી $g$ લાગુ કરીએ તોય
અનંત વખતનો જ અનુપ્રયોગ રહે છે. એટલે $Yg$ પર
$g$નો અનંત અનુપ્રયોગ જ મળે છે.

ધ્યાન આપો કે $Yg$થી શરૂ થતો ઉપરનો $\beta$-લઘુકરણક્રમ
અનંત છે. તેથી $Yg$ને કોઈ પદ પર લાગુ કરીને $(Yg)N$
વિચારીએ તો આ પદ પણ અનંત જુદાં પદોમાં લઘુકૃત થઈ શકે:
$(g(Yg))N$, $(g(g(Yg)))N$, \dots. તેમ છતાં લઘુકરણનો
કોઈ \emph{બીજો} ક્રમ સામાન્ય સ્વરૂપે સમાપ્ત થાય એવું શક્ય છે.

દાખલા તરીકે ક્રમગુણિત લો. $\fn{Fac}$ને $Y\,\fn{Fac}'$
તરીકે, એટલે $\fn{Fac'}$ના સ્થિરબિંદુ તરીકે, વ્યાખ્યાયિત
કરો. ત્યારે:
\begin{align*}
  \fn{Fac}\,\num 3
  &\red Y\, \fn{Fac}' \, \num 3 \\
  &\red \fn{Fac}' (Y \,\fn{Fac}') \, \num 3 \\
  &\ident (\lambd[x][\lambd[n][\fn{IsZero} \, n\, \num 1\,
      (\fn{Mult}\, n\, (x (\fn{Pred}\, n)))]]) \, \fn{Fac} \, \num 3\\
  & \red \fn{IsZero} \, \num 3\, \num 1\,
      (\fn{Mult}\, \num 3\, (\fn{Fac} (\fn{Pred}\, \num 3))) \\
  &\red  \fn{Mult}\, \num 3 \, (\fn{Fac} \, \num 2).
  \intertext{તે જ રીતે,}
  \fn{Fac}\,\num 2
  &\red  \fn{Mult}\, \num 2 \, (\fn{Fac} \, \num 1) \\
  \fn{Fac}\,\num 1
  &\red  \fn{Mult}\, \num 1 \, (\fn{Fac} \, \num 0)
  \intertext{પણ}
  \fn{Fac}\,\num 0
  &\red \fn{Fac}' (Y \,\fn{Fac}') \, \num 0 \\
  &\ident (\lambd[x][\lambd[n][\fn{IsZero} \, n\, \num 1\,
      (\fn{Mult}\, n\, (x (\fn{Pred}\, n)))]]) \, \fn{Fac} \, \num 0\\
  & \red \fn{IsZero} \, \num 0\, \num 1\,
      (\fn{Mult}\, \num 0\, (\fn{Fac} (\fn{Pred}\, \num 0))).\\
  &\red \num 1.
  \intertext{આથી એકંદરે}
  \fn{Fac}\,\num 3
  &\red  \fn{Mult}\, \num 3 \,
  (\fn{Mult} \, \num 2\, (\fn{Mult} \, \num 1\, \num 1)).
\end{align*}

$\fn{Fac'}$ માટે જે થયું તે કોઈ પણ પુનરાવર્તી વ્યાખ્યા
માટે થાય છે. ધારો કે આપણી પાસે નીચેનું પુનરાવર્તી
સમીકરણ છે:
\begin{align*}
g\,x_1\dots x_n & \equal[\beta] N
\intertext{અહીં $N$માં $g$ અને $x_1$, \dots,~$x_n$ આવી શકે છે. ત્યારે
હંમેશા એવું પદ $G \ident (Y \lambd[g][\lambd[x_1\dots x_n][N]])$
હોય છે કે}
G\,x_1 \dots x_n & \equal[\beta] \Subst{N}{G}{g}.
\intertext{કારણ કે સ્થિરબિંદુ પ્રમેયથી,}
G \ident (Y \lambd[g][\lambd[x_1\dots x_n][N]]) & \red \lambd[g][\lambd[x_1\dots x_n][N]](Y \lambd[g][\lambd[x_1\dots x_n][N]])\\
& \ident (\lambd[g][\lambd[x_1\dots x_n][N]])\,G
\intertext{અને તેથી}
G\,x_1 \dots x_n
& \red (\lambd[g][\lambd[x_1\dots x_n][N]])\,G\,x_1\dots x_n\\
& \red (\lambd[x_1\dots x_n][\Subst{N}{G}{g}])\,x_1\dots x_n\\
  & \red \Subst{N}{G}{g}.
\end{align*}

\olref{defn:Turing-Y}નો $Y$-સંયોજક એલન ટ્યુરિંગે આપ્યો હતો.
અલોન્ઝો ચર્ચે બીજો પ્રકાર સૂચવ્યો હતો, જેને આપણે~$Y_C$
કહીશું:
\[
Y_C \ident \lambd[g][(\lambd[x][g(xx)])(\lambd[x][g(xx)])].
\]
ચર્ચનો સંયોજક ટ્યુરિંગના સંયોજક કરતાં આ અર્થમાં થોડો
નબળો છે: $Y_Cg \equal[\beta] g(Y_Cg)$ છે, પરંતુ
$Y_Cg \bred g(Y_Cg)$ નથી.
\sourcecorrection{OLLAM-058}{સ્થિર અંગ્રેજી મૂળમાં
ચર્ચના સંયોજકના આ સરખામણીવાક્યમાં ભૂલથી
$Yg$ લખાયું છે. તે તો અગાઉના પ્રમેય મુજબ
$g(Yg)$માં લઘુકૃત થાય છે. અહીં ઉદ્દેશિત
$Y_Cg$ મૂક્યું છે.}
$V$ને $\lambd[x][g(xx)]$ પદ ગણો; ત્યારે
$Y_C \ident \lambd[g][VV]$. હવે
\begin{align*}
  VV & \ident (\lambd[x][g(xx)])V \red g(VV)
  \text{ અને તેથી}\\
  Y_C g & \ident (\lambd[g][VV])g \red VV \red g(VV), \text{ પરંતુ}\\
  g(Y_C g) & \ident g((\lambd[g][VV])g) \red g(VV).
\end{align*}
અન્ય શબ્દોમાં, $Y_Cg$ અને $g(Y_Cg)$ બંને એક જ પદ
$g(VV)$માં લઘુકૃત થાય છે; એટલે $Y_Cg \equal[\beta]
g(Y_Cg)$. ઉપયોગોમાં ઘણી વાર આટલું પૂરતું છે.

\end{document}
